% ===========================================================================
% MODELO DE RELATORIO  ·  ENE0046 Laboratorio de Eletronica  ·  2026/2
%
% Como usar
%   1. Preencha o bloco IDENTIFICACAO logo abaixo.
%   2. Deixe comentados os autores que nao existirem. Um autor fica
% centralizado,
%      dois dividem a linha, tres ocupam tercos.
%   3. Substitua o texto de cada secao pelo seu. As orientacoes que estao la
%      dizem o que se espera em cada uma.
%   4. Compile com pdflatex duas vezes, ou com latexmk -pdf.
%
% Estrutura
%   Roteiro 1 nao tem bancada: use so a subsecao 4.1 e a 5.1, e
% compare o calculo
%   analitico com a simulacao de projeto.
%   Roteiros 2 a 7 usam a estrutura inteira.
%
% O arquivo estilo-relatorio.tex e a imagem logo-unb.png devem estar na mesma
% pasta deste arquivo.
% ===========================================================================
\documentclass[11pt,a4paper]{article}
% !TeX root = modelo-relatorio.tex
% Formatacao do relatorio. Compile o modelo-relatorio.tex, nao este arquivo.
\usepackage[T1]{fontenc}
\usepackage[utf8]{inputenc}
\usepackage[brazil]{babel}
\usepackage[a4paper,top=1.6cm,bottom=1.8cm,left=2.1cm,right=2.1cm]{geometry}
\usepackage[default]{lato}
\usepackage{amsmath,amssymb}
\usepackage{siunitx}
\usepackage{graphicx}
\usepackage{booktabs,array,longtable,colortbl}
\usepackage{enumitem}
\usepackage{tcolorbox}
\usepackage{xcolor}
\usepackage{tikz}
\usepackage{titlesec}
\usepackage{fancyhdr}
\usepackage{caption}
\usepackage{etoolbox}
\usepackage[colorlinks=true,linkcolor=azul,urlcolor=azul,citecolor=azul]{hyperref}
\tcbuselibrary{skins,breakable}

\definecolor{azul}{HTML}{0A4B8C}
\definecolor{azulclaro}{HTML}{E8F0F8}
\definecolor{verde}{HTML}{0B7A4B}
\definecolor{cinza}{HTML}{5B6672}
\definecolor{cinzaclaro}{HTML}{F4F6F8}
\definecolor{vinho}{HTML}{B04A4A}
\definecolor{vinhoclaro}{HTML}{FBF2F2}

\sisetup{output-decimal-marker={,}, per-mode=symbol}
\setlength{\parindent}{0pt}
\setlength{\parskip}{6pt plus 1pt}
\setlist{itemsep=3pt, topsep=4pt, parsep=0pt, leftmargin=1.3em}
\renewcommand{\arraystretch}{1.25}
\linespread{1.06}

% ---- secoes numeradas com marcador colorido -------------------------------
\titleformat{\section}[hang]
  {\Large\bfseries\color{azul}}
  {\raisebox{-1pt}{\tikz\node[fill=azul,text=white,rounded corners=2pt,
     inner sep=3pt,minimum width=20pt,align=center]{\bfseries\thesection};}}
  {10pt}{}
\titlespacing*{\section}{0pt}{16pt plus 3pt}{6pt}
\titleformat{\subsection}{\large\bfseries\color{cinza}}{\thesubsection}{8pt}{}
\titlespacing*{\subsection}{0pt}{11pt plus 2pt}{3pt}

% ---- legendas -------------------------------------------------------------
\captionsetup{font=small,labelfont={bf,color=azul},labelsep=period,skip=6pt}

% ---- identificacao --------------------------------------------------------
\newcommand{\autorA}{}\newcommand{\matA}{}
\newcommand{\autorB}{}\newcommand{\matB}{}
\newcommand{\autorC}{}\newcommand{\matC}{}
\newcommand{\autorum}[2]{\renewcommand{\autorA}{#1}\renewcommand{\matA}{#2}}
\newcommand{\autordois}[2]{\renewcommand{\autorB}{#1}\renewcommand{\matB}{#2}}
\newcommand{\autortres}[2]{\renewcommand{\autorC}{#1}\renewcommand{\matC}{#2}}

\newcommand{\disciplina}{ENE0046 · Laboratório de Eletrônica}
\newcommand{\semestre}{2026/2}
\newcommand{\turma}{}
\newcommand{\roteironum}{}
\newcommand{\roteirotema}{}
\newcommand{\aturma}[1]{\renewcommand{\turma}{#1}}
\newcommand{\roteiro}[2]{\renewcommand{\roteironum}{#1}\renewcommand{\roteirotema}{#2}}

\newcommand{\umautor}[2]{%
  \begin{minipage}[t]{\dimexpr\linewidth-4pt\relax}\centering
    \parbox[t][30pt][t]{\linewidth}{\centering
      {\fontsize{10.5}{12.5}\selectfont #1}}\\[-2pt]
    {\color{cinza}\fontsize{9}{11}\selectfont\texttt{#2}}
  \end{minipage}}

\newcounter{nautores}
\newcommand{\blocoautores}{%
  \setcounter{nautores}{0}%
  \ifdefempty{\autorA}{}{\stepcounter{nautores}}%
  \ifdefempty{\autorB}{}{\stepcounter{nautores}}%
  \ifdefempty{\autorC}{}{\stepcounter{nautores}}%
  \ifnum\value{nautores}=0
    \begin{center}\color{vinho}\bfseries
      Preencha ao menos um autor no preâmbulo do arquivo.
    \end{center}
  \else
    \begin{center}
    \ifcase\value{nautores}\or
      % --- um autor: ocupa o centro
      \begin{minipage}{0.5\textwidth}\umautor{\autorA}{\matA}\end{minipage}%
    \or
      % --- dois autores: metades
      \begin{minipage}{0.42\textwidth}\umautor{\autorA}{\matA}\end{minipage}\hspace{0.04\textwidth}%
      \begin{minipage}{0.42\textwidth}\umautor{\autorB}{\matB}\end{minipage}%
    \or
      % --- tres autores: tercos
      \begin{minipage}{0.31\textwidth}\umautor{\autorA}{\matA}\end{minipage}\hspace{0.015\textwidth}%
      \begin{minipage}{0.31\textwidth}\umautor{\autorB}{\matB}\end{minipage}\hspace{0.015\textwidth}%
      \begin{minipage}{0.31\textwidth}\umautor{\autorC}{\matC}\end{minipage}%
    \fi
    \end{center}
  \fi}

% ---- cabecalho do relatorio -----------------------------------------------
\newcommand{\capa}{%
\begin{tcolorbox}[colback=azul,colframe=azul,boxrule=0pt,arc=4pt,
  left=16pt,right=16pt,top=13pt,bottom=13pt]
\color{white}
\begin{minipage}{0.76\textwidth}
  {\fontsize{9}{11}\selectfont\textbf{\MakeUppercase{\disciplina{} · \semestre{}}}}\\[6pt]
  {\fontsize{22}{25}\selectfont\bfseries Relatório do Roteiro \roteironum}\\[3pt]
  {\fontsize{12.5}{15}\selectfont \roteirotema}
\end{minipage}\hfill
\begin{minipage}{0.2\textwidth}\raggedleft
  \colorbox{white}{\hspace{2pt}\raisebox{-2pt}{\includegraphics[height=30pt]{logo-unb.png}}\hspace{2pt}}
\end{minipage}
\end{tcolorbox}
\vspace{10pt}
\blocoautores
\ifdefempty{\turma}{}{%
  \begin{center}{\color{cinza}\fontsize{9}{11}\selectfont Turma \turma}\end{center}}
\vspace{4pt}
{\color{azulclaro}\hrule height 2pt}
\vspace{2pt}}

% ---- caixas auxiliares ----------------------------------------------------
\newtcolorbox{formula}{colback=azulclaro,colframe=azul,boxrule=0.8pt,arc=3pt,
  left=12pt,right=12pt,top=9pt,bottom=9pt}
\newtcolorbox{nota}[1]{colback=cinzaclaro,colframe=cinza,boxrule=0.8pt,arc=3pt,
  left=10pt,right=10pt,top=7pt,bottom=7pt,
  title={\textbf{#1}},fonttitle=\small\bfseries,coltitle=white,colbacktitle=cinza}
\newenvironment{resumo}
  {\vspace{6pt}\begin{center}\begin{minipage}{0.92\textwidth}
   \begin{tcolorbox}[colback=cinzaclaro,colframe=cinzaclaro,boxrule=0pt,arc=3pt,
     left=14pt,right=14pt,top=10pt,bottom=10pt,
     borderline west={2.5pt}{0pt}{azul}]
   {\color{azul}\fontsize{8}{10}\selectfont\textbf{RESUMO}}\\[4pt]
   \small}
  {\end{tcolorbox}\end{minipage}\end{center}\vspace{2pt}}

\newtcolorbox{espec}{colback=azulclaro,colframe=azul,boxrule=0.8pt,arc=3pt,
  left=12pt,right=12pt,top=9pt,bottom=9pt,
  title={\textbf{Especificação da matrícula}},fonttitle=\bfseries,
  coltitle=white,colbacktitle=azul}

% ---- rodape ---------------------------------------------------------------
\pagestyle{fancy}\fancyhf{}
\renewcommand{\headrulewidth}{0pt}
\renewcommand{\footrulewidth}{0.4pt}
\fancyfoot[L]{\color{cinza}\fontsize{8}{10}\selectfont \disciplina{} · \semestre{} · Roteiro \roteironum}
\fancyfoot[R]{\color{cinza}\fontsize{8}{10}\selectfont \thepage}

\newlist{referencias}{itemize}{1}
\setlist[referencias]{label={}, leftmargin=1.1em, itemindent=-1.1em,
                     itemsep=3pt, topsep=2pt, parsep=0pt}

\newcommand{\mat}[1]{\texttt{\textbf{#1}}}


% --------------------------- IDENTIFICACAO ---------------------------------
\roteiro{4}{Amplificador Operacional: Nao-idealidades}
\aturma{01}

\autorum{Bernardo Manciola Lobo}{21/1026486}
\autordois{Enzo Morrone de Azevedo Muniz}{22/1035442}
\autortres{Matheus Vieira Tiveron}{24/1039879}
% ---------------------------------------------------------------------------

\begin{document}
\capa

\begin{resumo}
  Escreva este parágrafo por último, com 100 a 150 palavras. Diga o
  que você projetou, sob qual especificação, o que simulou ou mediu, o principal resultado
  numérico e a que
  conclusão chegou. Quem ler só este parágrafo deve saber o que você
  fez e o que encontrou.
\end{resumo}

\section{Introdução}

Comece pelo problema, e não pelo circuito. Que necessidade de
engenharia deu origem a esta
topologia, onde você a encontraria em um equipamento real e o que
aconteceria sem ela. Duas
ou três referências ajudam aqui.

Depois estreite o foco: que aspecto você vai investigar, sob quais
condições, e por que
vale a pena verificá-lo. Se o roteiro impõe alguma restrição de
projeto que parece
arbitrária, aproveite para explicar a razão dela.

No último parágrafo, diga qual é o objetivo deste relatório e como o texto está
organizado, seção a seção.

\section{Fundamentação teórica}

Deduza as expressões que você vai usar mais adiante, partindo das
leis de circuito.
Numere as equações que forem retomadas no restante do texto.

\begin{equation}
  H(j\omega) = \frac{1}{1+j\omega RC}
  \label{eq:exemplo}
\end{equation}

Toda grandeza que aparecer nos resultados precisa ter sido
apresentada aqui. Explique o que
cada termo significa fisicamente, não apenas o que a letra
representa. Registre as hipóteses
que você está assumindo, como linearidade ou idealidade dos
componentes: se alguma delas
não se sustentar no experimento, você vai retomá-la na discussão.

Se o circuito tiver regimes de funcionamento distintos, apresente
cada um e diga o que os
separa.

\section{Especificação e projeto}

Esta seção transcreve a especificação numérica de cada integrante,
calculada a partir dos nove dígitos da matrícula conforme as
Seções~2.2 a~2.5 do roteiro, e leva cada parâmetro até o valor de
componente na série E12. A Seção~3.1 reúne a especificação da equipe,
usada na bancada; a Seção~3.2 apresenta o projeto individual; e a
Seção~3.3 registra os valores medidos dos resistores efetivamente montados.

\subsection{Especificação da equipe}

A equipe~2 adotou, para a montagem em bancada, a especificação da
Tabela~\ref{tab:equipe}. Esses valores foram usados no projeto do
circuito único que a equipe montou, e são a referência contra a qual
a simulação da montagem e a medição são confrontadas.

\begin{table}[h]
  \centering
  \caption{Especificação da equipe~2, conforme a Seção~3.1 do roteiro.}
  \label{tab:equipe}
  \begin{tabular}{lc}
    \hline
    Grandeza & Valor \\
    $R_f$ e $R_g$ do ganho & $\SI{330}{\kilo\ohm}$ e $\SI{2,2}{\kilo\ohm}$ \\
    Passa baixas & $\SI{4,7}{\kilo\ohm}$ e $\SI{1}{\micro\farad}$ \\
    \hline
  \end{tabular}
\end{table}

Transcreva os parâmetros que a sua matrícula determina, mostrando a
conta a partir dos
dígitos.

\begin{espec}
  Matrícula \mat{00/0000000}. \\
  Preencha aqui os parâmetros calculados a partir dos dígitos,
  conforme o roteiro.
\end{espec}

Leve cada parâmetro até o valor de componente, deixando visível a
expressão que você
inverteu para chegar lá. Depois escolha o valor na série comercial e
calcule o quanto essa
escolha afasta você do valor ideal. A tabela abaixo organiza isso, e
é ela que depois
permite separar o que veio do projeto do que veio da montagem.

\begin{center}
  \begin{tabular}{@{}lrrr@{}}
    \toprule
    \textbf{Grandeza} & \textbf{Teórico} & \textbf{Adotado} & \textbf{Desvio} \\
    \midrule
    & & & \\
    \bottomrule
  \end{tabular}
\end{center}

Se mais de uma combinação atender à especificação, diga por que
escolheu esta, mostrando
que as outras ficavam mais distantes.

\section{Metodologia}

Conte o que você fez e com quais parâmetros, de modo que um colega
consiga repetir sem
precisar perguntar nada. Explique as escolhas que não são evidentes:
por que essa faixa de
varredura, por que esse passo de integração, por que esse intervalo de gravação.

Inclua a figura do esquemático, com legenda que descreva o que ela
mostra, e cite-a no
texto, como na Figura~\ref{fig:exemplo}. Se uma figura não for
mencionada em nenhum
parágrafo, ela provavelmente não precisa estar aqui.

\subsection{Simulação de projeto, uma por integrante}

teste

Cada integrante recebeu uma especificação diferente da própria
matrícula, e portanto chegou
a um dimensionamento diferente. Traga aqui \textbf{uma simulação por
integrante}, cada uma
com os componentes que aquela pessoa projetou. São elas que mostram
se cada projeto cumpre
o que foi pedido.

\subsection{Montagem em bancada - Elementos usados}

A equipe monta um circuito só. Diga \textbf{de quem} é a
especificação que vocês montaram e
por que escolheram essa. Informe os valores \textbf{medidos} de cada
componente que vocês
usaram, e não os valores nominais impressos no corpo da peça, porque
a diferença entre os
dois explica boa parte do que vocês vão observar.

\subsubsection{Amplificador de instrumentação}

Entrada foi $V_p = \SI{330,5}{\milli\volt}$

\begin{center}
  \begin{tabular}{@{}lrrrr@{}}
    \toprule
    \textbf{Grandeza} & \textbf{Nominal} & \textbf{Medido} & \textbf{Erro} \\
    \midrule
    $R_{gINA}$         & $\SI{10}{\kilo\ohm}$   & $\SI{9,892}{\kilo\ohm}$   & \\

    $R_{passa-altas}$  & $\SI{330}{\kilo\ohm}$  & $\SI{326,6}{\kilo\ohm}$   & \\
    $C_{passa-altas}$  & $\SI{1}{\micro\farad}$ & $\SI{1,03}{\micro\farad}$ & \\

    $R_{gGanho}$       & $\SI{330}{\kilo\ohm}$  & $\SI{326,7}{\kilo\ohm}$   & \\
    $R_{fGanho}$       & $\SI{2,2}{\kilo\ohm}$  & $\SI{2,177}{\kilo\ohm}$   & \\

    Ganho              & $\SI{151}{}$           & $\SI{150,7}{}$            & \\

    $R_{passa-baixas}$ & $\SI{4,7}{\kilo\ohm}$  & $\SI{4,632}{\kilo\ohm}$   & \\
    $C_{passa-baixas}$ & $\SI{1}{\micro\farad}$ & $\SI{1,03}{\micro\farad}$ & \\

    Divisor $R_1$      & $\SI{820}{\ohm}$       & $\SI{824}{\ohm}$          & \\
    Divisor $R_2$      & $\SI{100}{\kilo\ohm}$  & $\SI{100,8}{\kilo\ohm}$   & \\
    \bottomrule
  \end{tabular}
\end{center}

\subsection{Simulação da montagem}

Refaça a simulação usando os valores medidos dos componentes que
estavam na bancada. Ela
não substitui a simulação de projeto: é a comparação entre as duas
que separa o efeito da
tolerância dos componentes do restante.

\begin{nota}{Três comparações, três causas diferentes}
  Ao longo do relatório vocês vão confrontar três pares de
  resultados. Cada um responde a uma
  pergunta distinta, e vale mantê-los separados na hora de discutir.

  \smallskip
  \textbf{Cálculo analítico contra simulação de projeto.} A conta
  está certa? Se aqui houver
  diferença, procure na álgebra ou no esquemático.

  \smallskip
  \textbf{Simulação de projeto contra simulação da montagem.} Quanto
  custou a tolerância dos
  componentes? A diferença aqui é a distância entre o valor nominal e
  o valor real da peça
  que vocês pegaram.

  \smallskip
  \textbf{Simulação da montagem contra medição.} O modelo descreve o
  circuito real? Diferença
  aqui aponta algo que o modelo não representa, como a resistência
  série do indutor, a
  impedância da ponta de prova ou a banda do instrumento.
\end{nota}

\begin{figure}[htb]
  \centering
  % \includegraphics[width=0.8\linewidth]{arquivo.png}
  \fbox{\parbox[c][3.2cm][c]{0.8\linewidth}{\centering\color{cinza}
  substitua por \texttt{\textbackslash includegraphics}}}
  \caption{Legenda descritiva. Diga o que a figura mostra e o que
    nela é relevante, e não
  apenas o nome do circuito.}
  \label{fig:exemplo}
\end{figure}

\section{Resultados}

Apresente o que vocês obtiveram, em tabela sempre que houver mais de
dois números. Use o
número de algarismos significativos que a medição sustenta, e não
todos os que o programa
imprime. Para cada leitura de bancada, informe o instrumento e a incerteza.

Guarde a interpretação para a seção seguinte: aqui vocês mostram o
que encontraram.

\subsection{Projeto de cada integrante}

Uma tabela por integrante, ou uma tabela única com uma coluna por
integrante, confrontando
o valor analítico e o simulado para cada grandeza pedida no roteiro.

\subsection{Circuito montado}

Para cada grandeza, o que o projeto do integrante montado previa, o
que a simulação com os
valores medidos devolveu, e o que vocês mediram na bancada. É desta
tabela que sai a
discussão.

\subsubsection{Amplificador de instrumentação}
\begin{center}
  \begin{tabular}{@{}lrrrr@{}}
    \toprule
    \textbf{Grandeza} & \textbf{Projeto} & \textbf{Sim. montagem} &
    \textbf{Medido} & \textbf{Erro} \\
    \midrule
    Limite superior do ganho de modo comum & & & $\SI{0,6}{\milli\volt}$ & \\
    \bottomrule
  \end{tabular}
\end{center}

Em $\SI{0,6}{\hertz}$, a saída $V_p$ fica $0,707V_{10hz}$
\subsubsection{Passa altas}
\begin{center}
  \begin{tabular}{@{}lrrrr@{}}
    \toprule
    \textbf{Grandeza} & \textbf{Projeto} & \textbf{Sim. montagem} &
    \textbf{Medido} & \textbf{Erro} \\
    \midrule
    Frequência de corte                & $\SI{0,482}{\hertz}$ & &
    $\SI{0,6}{\hertz}$ & \\
    Saida $V_p$ em $\SI{10}{\hertz} $  &                      & &
    $\SI{2}{\volt}$    & \\
    Saida $V_p$ em $\SI{1}{\hertz}$    &                      & &
    $\SI{1,7}{\volt}$  & \\
    Saida $V_p$ em $\SI{0,8}{\hertz}$  &                      & &
    $\SI{1,6}{\volt}$  & \\
    Saida $V_p$ em $\SI{0,6}{\hertz}$  &                      & &
    $\SI{1,4}{\volt}$  & \\
    Saida $V_p$ em $\SI{0,45}{\hertz}$ &                      & &
    $\SI{1,18}{\volt}$ & \\
    \bottomrule
  \end{tabular}
\end{center}

\subsubsection{Ganho Adicional}
\begin{center}
  \begin{tabular}{@{}lrrrr@{}}
    \toprule
    \textbf{Grandeza} & \textbf{Projeto} & \textbf{Sim. montagem} &
    \textbf{Medido} &
    \textbf{Erro} \\
    \midrule
    Saída gerador & $\SI{1,35}{\volt}$ & & $\SI{1,35}{\volt}$ & \\
    Saída amplificador & $\SI{5}{\volt}$ & & $\SI{5}{\volt}$ & \\
    \bottomrule
  \end{tabular}
\end{center}

\subsubsection{Passa baixas}
\begin{center}
  \begin{tabular}{@{}lrrrr@{}}
    \toprule
    \textbf{Grandeza} & \textbf{Projeto} & \textbf{Sim. montagem} &
    \textbf{Medido} &
    \textbf{Erro} \\
    \midrule
    Saída gerador & $\SI{1,08}{\volt}$ & & $\SI{1,08}{\volt}$ & \\
    Saída estágio final & $\SI{2}{\volt}$ & & $\SI{2}{\volt}$ & \\

    Saída em $\SI{10}{\hertz} $ & $\SI{2}{\volt}$ & & $\SI{2}{\volt}$ & \\

    vpp
    Saída em $\SI{0,2}{\hertz} $ & $\SI{2}{\volt}$ & & $\SI{1,6}{\volt}$ & \\
    Saída em $\SI{3,5}{\hertz} $ & $\SI{2}{\volt}$ & & $\SI{3,84}{\volt}$ & \\
    Saída em $\SI{6,8}{\hertz} $ & $\SI{2}{\volt}$ & & $\SI{3,84}{\volt}$ & \\
    Saída em $\SI{10,1}{\hertz} $ & $\SI{2}{\volt}$ & & $\SI{3,76}{\volt}$ & \\
    Saída em $\SI{12,4}{\hertz} $ & $\SI{2}{\volt}$ & & $\SI{3,68}{\volt}$ & \\
    Saída em $\SI{15,7}{\hertz} $ & $\SI{2}{\volt}$ & & $\SI{3,56}{\volt}$ & \\
    Saída em $\SI{19,0}{\hertz} $ & $\SI{2}{\volt}$ & & $\SI{3,44}{\volt}$ & \\
    Saída em $\SI{22,3}{\hertz} $ & $\SI{2}{\volt}$ & & $\SI{3,28}{\volt}$ & \\
    Saída em $\SI{25,6}{\hertz} $ & $\SI{2}{\volt}$ & & $\SI{3,12}{\volt}$ & \\
    Saída em $\SI{28,9}{\hertz} $ & $\SI{2}{\volt}$ & & $\SI{2,96}{\volt}$ & \\
    Saída em $\SI{32,2}{\hertz} $ & $\SI{2}{\volt}$ & & $\SI{2,84}{\volt}$ & \\
    Saída em $\SI{35,5}{\hertz} $ & $\SI{2}{\volt}$ & & $\SI{2,72}{\volt}$ & \\
    Saída em $\SI{45,5}{\hertz} $ & $\SI{2}{\volt}$ & & $\SI{2,36}{\volt}$ & \\

    Frequência de corte & $\SI{33,8}{\hertz} $ & & $\SI{33,3}{\hertz} $ \\

    \bottomrule
  \end{tabular}
\end{center}

\subsubsection{Eletrocardiograma}
\begin{center}
  \begin{tabular}{@{}lrrrr@{}}
    \toprule
    \textbf{Grandeza} & \textbf{Projeto} & \textbf{Sim. montagem} &
    \textbf{Medido} &
    \textbf{Erro} \\
    \midrule
    Intervalo entre picos R & NA & NA & $\SI{784}{\milli\second}$ & \\

    \bottomrule
  \end{tabular}
\end{center}

\section{Discussão}

Percorra as três comparações na ordem em que aparecem no quadro da
Seção~4: a conta contra
a simulação de projeto, esta contra a simulação da montagem, e esta
contra a medição. Cada
passo isola uma causa, e a discussão fica bem mais curta quando eles
não se misturam.

Compare também os projetos dos integrantes entre si. Como as
especificações são diferentes,
os resultados diferem, e mostrar que essa diferença é a que a teoria
previa vale mais do
que repetir três vezes a mesma análise.

Sempre que uma divergência for maior do que a incerteza esperada,
aponte uma causa concreta
e diga como ela poderia ser reduzida. Vale, por exemplo, a tolerância
do componente, a
resistência série do indutor, a impedância de entrada do
osciloscópio, a banda do
amplificador ou a resolução da varredura.

Quando previsto e obtido coincidirem, explique por que era esperado
que coincidissem, e o
que essa coincidência não permite concluir.

Por fim, retome as hipóteses que vocês declararam na fundamentação e
diga quais se
confirmaram. Se algo observado contrariar a previsão, resolva a
contradição aqui.

\section{Conclusão}

Retome o objetivo que vocês anunciaram na introdução e responda se
ele foi atingido,
apoiando a resposta em um número. Aponte a principal limitação do que
foi feito e o que
ela impede de afirmar. Evite trazer resultado novo nesta seção.

\section{Referências}

\begin{referencias}
\item SEDRA, A.~S.; SMITH, K.~C. \emph{Microeletrônica}. 5.~ed. São
  Paulo: Pearson
  Prentice Hall, 2007.
\item BOYLESTAD, R.~L.; NASHELSKY, L. \emph{Dispositivos Eletrônicos e Teoria de
  Circuitos}. 11.~ed. São Paulo: Pearson Education do Brasil, 2013.
\end{referencias}

\end{document}
